\section{Introduction}


The exponential growth of blockchain state presents fundamental scalability challenges for decentralized networks. As of 2024, major blockchain platforms require hundreds of gigabytes of storage for full node operation, creating barriers to entry and centralizing network participation. This state growth problem manifests in three critical dimensions: storage requirements growing unbounded over time, synchronization bandwidth increasing linearly with state size, and verification costs scaling with tree depth.

Traditional Merkle Patricia trees, while providing cryptographic authentication of state, suffer from proof sizes that scale logarithmically with state size. For a blockchain with $2^{32}$ accounts, Merkle proofs require approximately 32 hash values, resulting in proofs exceeding 1 kilobyte. This overhead becomes prohibitive for light clients and cross-chain verification.

Verkle trees represent a paradigm shift in authenticated data structures by replacing hash-based commitments with polynomial commitments. This fundamental change enables constant-size proofs regardless of tree size or depth. A Verkle tree can prove membership or non-membership with a single 150-byte proof, compared to 640+ bytes for equivalent Merkle proofs.

The stateless client paradigm enabled by Verkle trees transforms blockchain architecture. Rather than maintaining complete state, clients can verify transactions using only block headers and compact proofs. This reduces storage requirements from hundreds of gigabytes to megabytes, enabling blockchain participation on resource-constrained devices including mobile phones and IoT devices.

Our contributions include:
\begin{itemize}
    \item Integration of Verkle trees with Lux's multi-consensus architecture
    \item Optimization of polynomial commitment schemes for parallel verification
    \item Migration framework preserving backward compatibility
    \item Performance analysis demonstrating practical feasibility
    \item Security analysis under adaptive adversarial models
\end{itemize}

